% !TEX spellcheck = en_US



% ~~~~~~~~~~~~~~~~~~~~~~~~~~~~~~~~~~~~~~~ V
% (NC) 2026 Haluk Bingol
% github.com/halukbingol/LaTeX-Templates
%
% Licensees may copy, distribute, display, and perform the work and make derivative works 
% and remixes based on it only for non-commercial purposes. 
% ~~~~~~~~~~~~~~~~~~~~~~~~~~~~~~~~~~~~~~~ A




% =====================================================================
%  Internationalization in Java -- Freshman Level: Your Program Meets the World
%  ---------------------------------------------------------------------
%  Built on Template/hblibTemplate.tex with the libraries in ../../hblib/.
%  Programs, terminal sessions and their real outputs are in codeoutput/:
%     codeoutput/<Name>.java  a program        codeoutput/<Name>.sh   a session
%     codeoutput/<Name>/      the session's project folder
%     codeoutput/<Name>.txt   what it printed, made by:  bash run_examples.sh
% =====================================================================

\documentclass[aspectratio=169,10pt,compress]{beamer}

% ~~~~~~~~~~~~~~~~~~~~~~~~~~~~~~~~~~~~~~~~~~~~~~~~~~~~~~~~~~~~~~~~~ V hblib
	\usepackage{../../hblib/hblibPresentation} % lib presentation: theme, i/N, \hSection, algorithm2e, ...
	\usepackage{../../hblib/hblibCodeoutput} % lib code + output: \lstcode, \lstterminal, ...
	\usepackage{../../hblib/hblibDatastructures} % lib data structures: \hString, \hStack, \hQueue
% ~~~~~~~~~~~~~~~~~~~~~~~~~~~~~~~~~~~~~~~~~~~~~~~~~~~~~~~~~~~~~~~~~ A hblib


% xcolor, array (and the L{width} column), listings, tikz and algorithm2e come from the hblib libraries
\usepackage[T1]{fontenc}
\usepackage{lmodern}
\usepackage{booktabs}

\definecolor{gitred}{RGB}{240,80,50}
\newcommand{\win}{\textbf{Windows}}
\newcommand{\mac}{\textbf{macOS}}

% algorithm2e set-up (the package is loaded by hblibPresentation), as in the template
\SetAlCapFnt{\small}
\SetAlCapNameFnt{\small}
\DontPrintSemicolon

% .properties files (resource bundles)
\lstdefinelanguage{properties}{morecomment=[l]{\#}, morecomment=[l]{!}}

% ---------------------------------------------------------------------
\title{Internationalization in Java}
\subtitle{Freshman Level: Your Program Meets the World}
\author[Bingol]{Haluk O. Bingol}

\institute[FCIS]{
    Faculty of Computer and Information Sciences\\
    Yeditepe University
}
\date{
	{\tiny \hbTimeStamp}
}

\begin{document}




% ~~~~~~~~~~~~~~~~~~~~~~~~~~~~~~~~~~~~~~ V frm
\begin{frame}[t,fragile]
  \titlepage
\end{frame}
% ~~~~~~~~~~~~~~~~~~~~~~~~~~~~~~~~~~~~~~ A frm




% ~~~~~~~~~~~~~~~~~~~~~~~~~~~~~~~~~~~~~~ V frm
\begin{frame}[t,fragile] \frametitle{Outline}
  \begin{columns}[T]
    \begin{column}{0.48\textwidth}
      \tableofcontents[sections={1-3}, hideallsubsections]
    \end{column}
    \begin{column}{0.48\textwidth}
      \tableofcontents[sections={4-20}, hideallsubsections]
    \end{column}
  \end{columns}
\end{frame}
% ~~~~~~~~~~~~~~~~~~~~~~~~~~~~~~~~~~~~~~ A frm




% ~~~~~~~~~~~~~~~~~~~~~~~~~~~~~~~~~~~~~~ V frm
\begin{frame}[t,fragile] \frametitle{Your program will run on other computers}
  \begin{columns}[T]
    \begin{column}{0.52\textwidth}
      \begin{itemize}\small
        \item Your laptop, your friend's laptop, the lab, the teacher's grading script, a server abroad:
          each may have a \textbf{different language and region}.
        \item Java then prints and reads numbers differently, sorts and capitalizes differently,
          and may show Turkish letters as \texttt{?}.
        \item This deck: the \textbf{first-week surprises} and the \textbf{simple fixes}.
      \end{itemize}
    \end{column}
    \begin{column}{0.44\textwidth}
      \begin{block}{Try it on your computer}
        \small
        \win: \menu{Settings > Time \& language > Language \& region}\\[0.3em]
        \mac: \menu{System Settings > General > Language \& Region}\\[0.3em]
        Look at \emph{Regional format} / \emph{Region}: this is what Java uses.
      \end{block}
      \small On these slides \texttt{LANG=tr\_TR.UTF-8} in front of a command plays a Turkish computer
      (macOS and Linux; on \win{} change the setting above).
    \end{column}
  \end{columns}
\end{frame}
% ~~~~~~~~~~~~~~~~~~~~~~~~~~~~~~~~~~~~~~ A frm

% ~~~~~~~~~~~~~~~~~~~~~~~~~~~~~~~~~~~~~~ V
\hSection[Printing]{Printing Numbers}
% ~~~~~~~~~~~~~~~~~~~~~~~~~~~~~~~~~~~~~~ A




% ~~~~~~~~~~~~~~~~~~~~~~~~~~~~~~~~~~~~~~ V frm
\begin{frame}[t,fragile] \frametitle{Same program, two computers}
  \begin{columns}[T]
    \begin{column}{0.5\textwidth}
      \lstcode[basicstyle=\ttfamily\scriptsize]{codeoutput/Receipt/Receipt.java}{Java}
    \end{column}
    \begin{column}{0.48\textwidth}
      \uncover<2->{\lstterminal[basicstyle=\ttfamily\scriptsize]{codeoutput/Receipt.txt}}
    \end{column}
  \end{columns}
  \vspace{0.1em}
  \small \texttt{printf} and \texttt{String.format} use the computer's locale: \texttt{1249,90} with a
  \textbf{comma} in Turkish, and \texttt{,} groups thousands with a \textbf{point}.
  Plain \texttt{+} always prints Java's own way (\texttt{3749.7000000000003}).
\end{frame}
% ~~~~~~~~~~~~~~~~~~~~~~~~~~~~~~~~~~~~~~ A frm




% ~~~~~~~~~~~~~~~~~~~~~~~~~~~~~~~~~~~~~~ V frm
\begin{frame}[t,fragile] \frametitle{Choose on purpose}
  \begin{columns}[T]
    \begin{column}{0.56\textwidth}
      \lstcode[basicstyle=\ttfamily\tiny]{codeoutput/ReceiptFixed/Receipt.java}{Java}
    \end{column}
    \begin{column}{0.42\textwidth}
      \uncover<2->{\lstterminal[basicstyle=\ttfamily\tiny]{codeoutput/ReceiptFixed.txt}}
    \end{column}
  \end{columns}
  \vspace{0.2em}
  \begin{itemize}\small
    \item Output for \textbf{programs} (files, CSV, a grading script): give \texttt{Locale.US} or \texttt{Locale.ROOT}.
    \item Output for \textbf{people}: the user's locale, or a locale you choose on purpose.
  \end{itemize}
\end{frame}
% ~~~~~~~~~~~~~~~~~~~~~~~~~~~~~~~~~~~~~~ A frm

% ~~~~~~~~~~~~~~~~~~~~~~~~~~~~~~~~~~~~~~ V
\hSection[Reading]{Reading Numbers}
% ~~~~~~~~~~~~~~~~~~~~~~~~~~~~~~~~~~~~~~ A




% ~~~~~~~~~~~~~~~~~~~~~~~~~~~~~~~~~~~~~~ V frm
\begin{frame}[t,fragile] \frametitle{\texttt{Scanner} reads the local way, too}
  \begin{columns}[T]
    \begin{column}{0.46\textwidth}
      \lstcode[basicstyle=\ttfamily\scriptsize]{codeoutput/Average/Average.java}{Java}
    \end{column}
    \begin{column}{0.52\textwidth}
      \uncover<2->{\lstterminal[basicstyle=\ttfamily\tiny]{codeoutput/Average.txt}}
    \end{column}
  \end{columns}
  \vspace{0.1em}
  \small On a Turkish computer \texttt{nextDouble()} wants \texttt{3,5}; \texttt{3.5} gives
  \texttt{InputMismatchException}. The program is right, the \textbf{input format} is wrong for that locale.
\end{frame}
% ~~~~~~~~~~~~~~~~~~~~~~~~~~~~~~~~~~~~~~ A frm




% ~~~~~~~~~~~~~~~~~~~~~~~~~~~~~~~~~~~~~~ V frm
\begin{frame}[t,fragile] \frametitle{Fixing the input format}
  \begin{columns}[T]
    \begin{column}{0.6\textwidth}
      \lstcode[basicstyle=\ttfamily\tiny]{codeoutput/AverageFixed/Average.java}{Java}
    \end{column}
    \begin{column}{0.38\textwidth}
      \uncover<2->{\lstterminal[basicstyle=\ttfamily\tiny]{codeoutput/AverageFixed.txt}}
    \end{column}
  \end{columns}
  \vspace{0.2em}
  \small \texttt{useLocale(Locale.ROOT)}: input with a decimal \textbf{point} on every computer.
  Use it for test data and files; keep the default when people type numbers the way they are used to.
\end{frame}
% ~~~~~~~~~~~~~~~~~~~~~~~~~~~~~~~~~~~~~~ A frm




% ~~~~~~~~~~~~~~~~~~~~~~~~~~~~~~~~~~~~~~ V frm
\begin{frame}[t,fragile] \frametitle{\texttt{parseDouble} knows only the point}
  \begin{columns}[T]
    \begin{column}{0.54\textwidth}
      \lstcode[basicstyle=\ttfamily\scriptsize]{codeoutput/ParseError/Grade.java}{Java}
    \end{column}
    \begin{column}{0.44\textwidth}
      \uncover<2->{\lstterminal[basicstyle=\ttfamily\tiny]{codeoutput/ParseError.txt}}
    \end{column}
  \end{columns}

  \vspace{0.1em}
  \small \texttt{Double.parseDouble} and \texttt{Integer.parseInt} never look at the locale:
  \texttt{"85,5"} is always an error. To accept \texttt{85,5}, use \texttt{NumberFormat.getInstance(locale).parse(s)}
  (next decks).
\end{frame}
% ~~~~~~~~~~~~~~~~~~~~~~~~~~~~~~~~~~~~~~ A frm

% ~~~~~~~~~~~~~~~~~~~~~~~~~~~~~~~~~~~~~~ V
\hSection[Letters]{Turkish Letters}
% ~~~~~~~~~~~~~~~~~~~~~~~~~~~~~~~~~~~~~~ A




% ~~~~~~~~~~~~~~~~~~~~~~~~~~~~~~~~~~~~~~ V frm
\begin{frame}[t,fragile] \frametitle{Where do the \texttt{?} come from?}
  \begin{columns}[T]
    \begin{column}{0.42\textwidth}
      \lstcode[basicstyle=\ttfamily\scriptsize]{codeoutput/Letters/Letters.java}{Java}
    \end{column}
    \begin{column}{0.56\textwidth}
      \uncover<2->{\lstterminal[basicstyle=\ttfamily\tiny]{codeoutput/Letters.txt}}
    \end{column}
  \end{columns}

  \vspace{0.1em}
  \small The \texttt{String} is fine; the \textbf{bytes} sent to the console depend on \texttt{stdout.encoding}.
  US-ASCII has no \texttt{ç}: \texttt{?}. UTF-8 uses 2 bytes for \texttt{Ç} (\texttt{c3 87}), ISO-8859-9 one (\texttt{c7}):
  the console must expect the same encoding.
\end{frame}
% ~~~~~~~~~~~~~~~~~~~~~~~~~~~~~~~~~~~~~~ A frm




% ~~~~~~~~~~~~~~~~~~~~~~~~~~~~~~~~~~~~~~ V frm
\begin{frame}[t,fragile] \frametitle{Turkish letters on Windows and macOS}
  \begin{columns}[T]
    \begin{column}{0.48\textwidth}
      \begin{block}{Windows (Command Prompt)}


% +++++++++++++++++++++++++++++++++++++++ V lst
%: -Lst.
\begin{lstlisting}[style=codesnippet, language={}, basicstyle=\ttfamily\scriptsize]
chcp 65001
java Letters.java
rem if letters are still wrong:
java -Dstdout.encoding=UTF-8 Letters.java
\end{lstlisting}
% +++++++++++++++++++++++++++++++++++++++ A lst


      \end{block}
    \end{column}
    \begin{column}{0.48\textwidth}
      \begin{block}{macOS (Terminal)}


% +++++++++++++++++++++++++++++++++++++++ V lst
%: -Lst.
\begin{lstlisting}[style=codesnippet, language=bash, basicstyle=\ttfamily\scriptsize]
java Letters.java
# Terminal uses UTF-8: nothing to do
locale          # shows LANG, LC_ALL, ...
\end{lstlisting}
% +++++++++++++++++++++++++++++++++++++++ A lst


      \end{block}
    \end{column}
  \end{columns}

  \vspace{0.3em}
  \begin{itemize}\small
    \item \win{} Command Prompt uses an old \emph{code page} (857 on Turkish Windows, 437 on English Windows);
      \texttt{chcp 65001} switches it to UTF-8 for this window. Windows Terminal and the IDE consoles are usually fine.
    \item Save your \texttt{.java} files as \textbf{UTF-8} (all editors can). Since Java 18 \texttt{javac} reads UTF-8 by default;
      for older Java use \texttt{javac -encoding UTF-8}.
  \end{itemize}
\end{frame}
% ~~~~~~~~~~~~~~~~~~~~~~~~~~~~~~~~~~~~~~ A frm

% ~~~~~~~~~~~~~~~~~~~~~~~~~~~~~~~~~~~~~~ V
\hSection[Case]{Upper and Lower Case}
% ~~~~~~~~~~~~~~~~~~~~~~~~~~~~~~~~~~~~~~ A




% ~~~~~~~~~~~~~~~~~~~~~~~~~~~~~~~~~~~~~~ V frm
\begin{frame}[t,fragile] \frametitle{The bug that only Turkish users see}
  \begin{columns}[T]
    \begin{column}{0.54\textwidth}
      \lstcode[basicstyle=\ttfamily\tiny]{codeoutput/Menu/Menu.java}{Java}
    \end{column}
    \begin{column}{0.44\textwidth}
      \uncover<2->{\lstterminal[basicstyle=\ttfamily\scriptsize]{codeoutput/Menu.txt}}
    \end{column}
  \end{columns}
  \vspace{0.1em}
  \small In Turkish, upper case \texttt{i} is \texttt{İ}, so \texttt{"quit"} becomes \texttt{"QUİT"},
  which is not \texttt{"QUIT"}. The program works in the lab and fails for the user.
\end{frame}
% ~~~~~~~~~~~~~~~~~~~~~~~~~~~~~~~~~~~~~~ A frm




% ~~~~~~~~~~~~~~~~~~~~~~~~~~~~~~~~~~~~~~ V frm
\begin{frame}[t,fragile] \frametitle{The fix: \texttt{Locale.ROOT}}
  \begin{columns}[T]
    \begin{column}{0.6\textwidth}
      \lstcode[basicstyle=\ttfamily\tiny]{codeoutput/MenuFixed/Menu.java}{Java}
    \end{column}
    \begin{column}{0.38\textwidth}
      \uncover<2->{\lstterminal[basicstyle=\ttfamily\scriptsize]{codeoutput/MenuFixed.txt}}
    \end{column}
  \end{columns}

  \vspace{0.1em}
  \small Commands, keywords, file extensions: \texttt{toUpperCase(Locale.ROOT)} or
  \texttt{equalsIgnoreCase}. Names and sentences shown to people: \texttt{toUpperCase(userLocale)}.
\end{frame}
% ~~~~~~~~~~~~~~~~~~~~~~~~~~~~~~~~~~~~~~ A frm

% ~~~~~~~~~~~~~~~~~~~~~~~~~~~~~~~~~~~~~~ V
\hSection[Translation]{A First Translation}
% ~~~~~~~~~~~~~~~~~~~~~~~~~~~~~~~~~~~~~~ A




% ~~~~~~~~~~~~~~~~~~~~~~~~~~~~~~~~~~~~~~ V frm
\begin{frame}[t,fragile] \frametitle{Texts in files, not in code}
  \begin{columns}[T]
    \begin{column}{0.56\textwidth}
      \lstcode[basicstyle=\ttfamily\tiny]{codeoutput/Greeter/Greeter.java}{Java}
    \end{column}
    \begin{column}{0.42\textwidth}
      \lstcode[basicstyle=\ttfamily\tiny, numbers=none]{codeoutput/Greeter/Texts.properties}{properties}
      \lstcode[basicstyle=\ttfamily\tiny, numbers=none]{codeoutput/Greeter/Texts_tr.properties}{properties}
    \end{column}
  \end{columns}
  \vspace{0.1em}
  \small \texttt{getBundle("Texts")} looks for \texttt{Texts\_tr.properties} on a Turkish computer and
  \texttt{Texts.properties} otherwise. Files go next to the classes (on the class path).
\end{frame}
% ~~~~~~~~~~~~~~~~~~~~~~~~~~~~~~~~~~~~~~ A frm




% ~~~~~~~~~~~~~~~~~~~~~~~~~~~~~~~~~~~~~~ V frm
\begin{frame}[t,fragile] \frametitle{One program, two languages}
  \lstterminal[basicstyle=\ttfamily\scriptsize]{codeoutput/Greeter.txt}

  \vspace{0.1em}
  \small
  No \texttt{if (turkish)} in the code: adding French means adding \texttt{Texts\_fr.properties}.
\end{frame}
% ~~~~~~~~~~~~~~~~~~~~~~~~~~~~~~~~~~~~~~ A frm




% ~~~~~~~~~~~~~~~~~~~~~~~~~~~~~~~~~~~~~~ V frm
\begin{frame}[t,fragile] \frametitle{First errors with bundles}
  \lstterminal[basicstyle=\ttfamily\tiny]{codeoutput/Missing.txt}

  \vspace{0.1em}
  \small
  1: run from another folder: the \texttt{.properties} files are not on the class path, \texttt{MissingResourceException}
  for the \emph{bundle}. 2: a typo in the Turkish file: the key is taken from the base file (English).
  3: a key that exists nowhere: \texttt{MissingResourceException} for the \emph{key}.
\end{frame}
% ~~~~~~~~~~~~~~~~~~~~~~~~~~~~~~~~~~~~~~ A frm

% ~~~~~~~~~~~~~~~~~~~~~~~~~~~~~~~~~~~~~~ V
\hSection[Summary]{Summary}
% ~~~~~~~~~~~~~~~~~~~~~~~~~~~~~~~~~~~~~~ A




% ~~~~~~~~~~~~~~~~~~~~~~~~~~~~~~~~~~~~~~ V frm
\begin{frame}[t,fragile] \frametitle{Cheat sheet}
  \begin{center}\small
  \begin{tabular}{L{0.36\textwidth}L{0.56\textwidth}}
    \toprule
    \textbf{problem} & \textbf{fix} \\
    \midrule
    \texttt{3,14} instead of \texttt{3.14}      & \texttt{printf(Locale.US, "\%.2f", x)} for programs \\
    \texttt{InputMismatchException}             & \texttt{new Scanner(System.in).useLocale(Locale.ROOT)} \\
    \texttt{"85,5"} gives \texttt{NumberFormatException} & \texttt{NumberFormat.getInstance(locale).parse(s)} \\
    \texttt{?} instead of \texttt{ç}            & \win: \texttt{chcp 65001}; files saved as UTF-8 \\
    \texttt{"quit".toUpperCase()} is \texttt{"QUİT"} & \texttt{toUpperCase(Locale.ROOT)} \\
    texts in two languages                      & \texttt{ResourceBundle.getBundle("Texts")} + \texttt{.properties} files \\
    \texttt{MissingResourceException}           & files on the class path? key spelled right? \\
    \bottomrule
  \end{tabular}
  \end{center}
\end{frame}
% ~~~~~~~~~~~~~~~~~~~~~~~~~~~~~~~~~~~~~~ A frm




% ~~~~~~~~~~~~~~~~~~~~~~~~~~~~~~~~~~~~~~ V frm
\begin{frame}[t,fragile] \frametitle{Do and do not}
  \begin{columns}[T]
    \begin{column}{0.48\textwidth}
      \begin{block}{Do}
        \small
        \begin{itemize}
          \item test your program once with a Turkish locale
          \item use \texttt{Locale.ROOT} for keywords and data files
          \item save sources as UTF-8
          \item keep user texts in \texttt{.properties} files
        \end{itemize}
      \end{block}
    \end{column}
    \begin{column}{0.48\textwidth}
      \begin{alertblock}{Do not}
        \small
        \begin{itemize}
          \item assume everyone writes \texttt{3.5}
          \item compare commands after \texttt{toUpperCase()} without a locale
          \item glue sentences: \texttt{"Hello, " + name + "!"} cannot be translated well
          \item ignore \texttt{?} in the output: the data may already be lost
        \end{itemize}
      \end{alertblock}
    \end{column}
  \end{columns}
\end{frame}
% ~~~~~~~~~~~~~~~~~~~~~~~~~~~~~~~~~~~~~~ A frm




% ~~~~~~~~~~~~~~~~~~~~~~~~~~~~~~~~~~~~~~ V frm
\begin{frame}[t,fragile] \frametitle{Exercises}
  \begin{enumerate}\small
    \item Change your regional format to Turkish, run \texttt{Receipt}, then change it back.
    \item Write a program that reads prices in Turkish style (\texttt{12,50}) and prints the total in US style.
    \item Find the bug: \texttt{if (ext.toLowerCase().equals("ini"))} for the file \texttt{SETUP.INI}. On which computers does it fail? Why?
    \item Add \texttt{Texts\_de.properties} to the greeter and run it with \texttt{-Duser.language=de}.
    \item Make the greeter print the date in the user's language (look at \texttt{DateTimeFormatter.ofLocalizedDate}).
  \end{enumerate}
\end{frame}
% ~~~~~~~~~~~~~~~~~~~~~~~~~~~~~~~~~~~~~~ A frm




% ~~~~~~~~~~~~~~~~~~~~~~~~~~~~~~~~~~~~~~ V frm
\begin{frame}[t,fragile]
	\begin{center}
		\vspace{4cm}
		\hRemark{\Huge Questions?}
	\end{center}
\end{frame}
% ~~~~~~~~~~~~~~~~~~~~~~~~~~~~~~~~~~~~~~ A frm



\end{document}
